\documentclass[12pt]{article}
\usepackage[utf8]{inputenc}
\usepackage[T1]{fontenc}
\usepackage{hyperref}
\title{Southwestern Tepehuan Language Resources for K-12 Curriculum Development}
\author{Compiled by Claude Sonnet 5 (claude-sonnet-5) for Andr\'as Kornai}
\date{}
\begin{document}
\maketitle

\section{Introduction}
Southwestern Tepehuan (ISO 639-3 tla; Audam, O'dam) is a Southern Tepehuan variety of the Uto-Aztecan family, spoken in the highlands of southwestern Durango, Mexico. UNESCO classes it as vulnerable. Per the search results, Audam has very little published grammatical description. The better-documented variety is Southeastern Tepehuan, which is a different language and should not be conflated with it.

\section{Primary Educational Resources}
None found for Southwestern Tepehuan.

\section{Linguistic and Reference Materials}
\begin{itemize}
\item Language snapshot of Southern Tepehuan, Garc\'ia Salido \& Everdell (2020), \emph{Language Documentation and Description} 19: \url{https://www.lddjournal.org/article/id/1243}
\item Tepehu\'an language overview (Wikipedia): \url{https://en.wikipedia.org/wiki/Tepehu%C3%A1n_language}
\item Endangered Languages Project entry: \url{https://endangeredlanguages.com/node/112162}
\item For the Southeastern variety only (not tla): Webonary grammar pages, \url{https://www.webonary.work/tepehuano-sureste/language/grammar/?lang=en}. The Willett \& Willett (2015) SIL dictionary of Santa Mar\'ia Ocotán was mentioned in results, but no URL was seen for it.
\end{itemize}

\section{Religious and Community Resources}
No Southwestern Tepehuan Bible translation or community organization turned up.

\section{Digital and Audio Resources}
Nothing specific to Southwestern Tepehuan was found.

\section{Curriculum Development Context}
Southeastern Tepehuan materials might serve as a comparison, but they are a different variety. Materials for tla would need new development with the Audam community.
\end{document}